% 第二章：随机变量
\chapter{随机变量}

\section{随机变量及其分布函数}
\label{sec:random-variables-and-distribution-functions}

\subsection*{随机变量的定义}

\begin{mydefinition}{随机变量}{random-variable}
给定概率空间 $(\Omega,\mathscr F,P)$。若映射 $X:\Omega\to\mathbb R$ 满足：对任意 $B\in\mathscr B(\mathbb R)$，都有
\[
X^{-1}(B):=\{\omega\in\Omega:X(\omega)\in B\}\in\mathscr F,
\]
则称 $X$ 为该概率空间上的\textbf{随机变量}。随机变量通常用 $X,Y,Z$ 或 $\xi,\eta$ 等表示。
\end{mydefinition}

\textbf{注：}\quad $X^{-1}(B)$ 表示集合 $B$ 的原像，不要求 $X$ 有反函数。随机变量在每个给定样本点上的值是确定的；其取值的随机性来自试验结果的随机性。

\begin{proposition}{随机变量的半直线判别}{random-variable-halfline-criterion}
映射 $X:\Omega\to\mathbb R$ 是随机变量，当且仅当对任意 $t\in\mathbb R$ 都有
\[
\{\omega\in\Omega:X(\omega)\leq t\}\in\mathscr F.
\]
\end{proposition}

\begin{example}[示性函数]
设 $A\in\mathscr F$，定义
\[
I_A(\omega)=
\begin{cases}
1,&\omega\in A,\\
0,&\omega\notin A.
\end{cases}
\]
称 $I_A$ 为事件 $A$ 的\textbf{示性函数}。由于
\[
\{I_A\leq t\}=
\begin{cases}
\varnothing,&t<0,\\
A^{\mathrm c},&0\leq t<1,\\
\Omega,&t\geq1,
\end{cases}
\]
上述集合均属于 $\mathscr F$，所以 $I_A$ 是随机变量。并且
\[
P(I_A=1)=P(A),\qquad P(I_A=0)=P(A^{\mathrm c}).
\]
\end{example}

\textbf{注：}\quad 为简化记号，通常写
\[
\{X\in B\}=X^{-1}(B)=\{\omega\in\Omega:X(\omega)\in B\},
\]
并将其概率记为
\[
P(X\in B)=P\bigl(\{\omega\in\Omega:X(\omega)\in B\}\bigr).
\]
例如，
\[
\{X=a\}=\{\omega:X(\omega)=a\},
\qquad
\{a<X\leq b\}=\{\omega:a<X(\omega)\leq b\}.
\]

\subsection*{分布函数}

对任意 $a<b$，有
\[
P(a<X\leq b)=P(X\leq b)-P(X\leq a).
\]
因此，知道所有 $P(X\leq t)$，就可以直接计算半开区间上的概率。

\begin{mydefinition}{分布函数}{distribution-function}
设 $X$ 是概率空间 $(\Omega,\mathscr F,P)$ 上的随机变量。函数
\[
F(t):=P(X\leq t),\qquad t\in\mathbb R,
\]
称为随机变量 $X$ 的\textbf{分布函数}。需要区分不同随机变量时，也记为 $F_X(t)$。
\end{mydefinition}

\textbf{注：}\quad 分布函数的定义域为 $\mathbb R$，且 $0\leq F(t)\leq1$。其实际值域是 $[0,1]$ 的子集，不必取遍 $[0,1]$。

\subsection*{分布函数的性质}

\begin{mytheorem}{分布函数的基本性质}{distribution-function-properties}
设 $F$ 为随机变量 $X$ 的分布函数，则：
\begin{enumerate}
    \item $F$ 单调不减，即 $a\leq b$ 时有 $F(a)\leq F(b)$；
    \item $\displaystyle\lim_{t\to-\infty}F(t)=0$，$\displaystyle\lim_{t\to+\infty}F(t)=1$；
    \item $F$ 右连续，即对任意 $a\in\mathbb R$，
    \[
    F(a)=\lim_{t\to a^+}F(t);
    \]
    \item 对任意 $a\in\mathbb R$，记
    \[
    F(a^-):=\lim_{t\to a^-}F(t),
    \]
    则
    \[
    P(X=a)=F(a)-F(a^-).
    \]
\end{enumerate}
\end{mytheorem}

\begin{proof}
（1）由 $\{X\leq a\}\subseteq\{X\leq b\}$ 和概率的单调性，得到 $F(a)\leq F(b)$。

（2）由于每个样本点上的 $X(\omega)$ 都是有限实数，
\[
\{X\leq-n\}\downarrow\varnothing,
\qquad
\{X\leq n\}\uparrow\Omega.
\]
由概率的上、下连续性，
\[
\lim_{n\to\infty}F(-n)=0,
\qquad
\lim_{n\to\infty}F(n)=1.
\]
结合 $F$ 的单调性，便得到 $t\to-\infty$ 和 $t\to+\infty$ 时的结论。

（3）取任意从右侧趋于 $a$ 的数列 $t_n\downarrow a$，则
\[
\{X\leq t_n\}\downarrow\{X\leq a\}.
\]
由概率的上连续性，
\[
\lim_{n\to\infty}F(t_n)=P(X\leq a)=F(a),
\]
所以 $F$ 右连续。

（4）取 $t_n<a$ 且 $t_n\uparrow a$，则
\[
\{X\leq t_n\}\uparrow\{X<a\}.
\]
由概率的下连续性及单调函数左极限的存在性，
\[
F(a^-)=P(X<a).
\]
于是
\[
P(X=a)=P(X\leq a)-P(X<a)=F(a)-F(a^-).
\]
\end{proof}

\textbf{注：}\quad $F$ 在 $a$ 处的跳跃大小恰好等于 $P(X=a)$。因此，$F$ 在 $a$ 处连续当且仅当 $P(X=a)=0$。

\textbf{注：}\quad 单调不减、右连续以及在负无穷和正无穷处分别趋于 $0$、$1$，也是分布函数的充分条件：任何满足这三个条件的函数，都可以成为某个概率空间上某个随机变量的分布函数。

\subsection*{用分布函数计算概率}

利用 $F(a^-)=P(X<a)$，对 $a<b$ 可得
\[
\begin{aligned}
P(a<X\leq b)&=F(b)-F(a),\\
P(a\leq X\leq b)&=F(b)-F(a^-),\\
P(a<X<b)&=F(b^-)-F(a),\\
P(a\leq X<b)&=F(b^-)-F(a^-).
\end{aligned}
\]
此外，
\[
P(X>a)=1-F(a),
\qquad
P(X\geq a)=1-F(a^-).
\]
计算时应先检查区间端点是否包含，以及端点处是否有正概率。

\begin{example}[含跳跃的分布函数]
随机变量 $X$ 的分布函数为
\[
F(t)=
\begin{cases}
0,&t<0,\\[0.3em]
\dfrac t2,&0\leq t<1,\\[0.5em]
\dfrac23,&1\leq t<2,\\[0.5em]
\dfrac{11}{12},&2\leq t<3,\\[0.5em]
1,&t\geq3.
\end{cases}
\]
分别计算：
\[
P(X<3)=F(3^-)=\frac{11}{12},
\]
\[
P(X=1)=F(1)-F(1^-)=\frac23-\frac12=\frac16,
\]
\[
P\left(X>\frac12\right)=1-F\left(\frac12\right)
=1-\frac14=\frac34,
\]
\[
P(2<X\leq4)=F(4)-F(2)=1-\frac{11}{12}=\frac1{12},
\]
\[
P\left(\frac12\leq X<2\right)
=F(2^-)-F\left(\left(\frac12\right)^-\right)
=\frac23-\frac14=\frac5{12}.
\]
\end{example}

\subsection*{分布函数的凸组合}

\begin{example}[可列个分布函数的凸组合]
设 $\{F_m\}_{m\geq1}$ 是一列分布函数，$a_m\geq0$ 且
\[
\sum_{m=1}^{\infty}a_m=1.
\]
令
\[
F(t)=\sum_{m=1}^{\infty}a_mF_m(t).
\]
称 $F$ 为这些分布函数的\textbf{凸组合}。证明 $F$ 也是分布函数。

首先，
\[
0\leq\sum_{m=1}^{\infty}a_mF_m(t)
\leq\sum_{m=1}^{\infty}a_m=1,
\]
所以级数收敛，且 $F$ 定义良好。

（1）若 $a\leq b$，则每个 $F_m(a)\leq F_m(b)$，从而
\[
F(a)=\sum_{m=1}^{\infty}a_mF_m(a)
\leq\sum_{m=1}^{\infty}a_mF_m(b)=F(b).
\]

（2）固定 $a\in\mathbb R$ 和 $\varepsilon>0$。选择正整数 $M$，使
\[
\sum_{m=M+1}^{\infty}a_m<\frac{\varepsilon}{2}.
\]
由 $F_1,\ldots,F_M$ 右连续，存在 $\delta>0$，使 $a<t<a+\delta$ 时，对所有 $1\leq m\leq M$ 都有
\[
0\leq F_m(t)-F_m(a)<\frac{\varepsilon}{2}.
\]
因此
\[
\begin{aligned}
0\leq F(t)-F(a)
&=\sum_{m=1}^{M}a_m\bigl(F_m(t)-F_m(a)\bigr)\\
&\quad+\sum_{m=M+1}^{\infty}a_m\bigl(F_m(t)-F_m(a)\bigr)\\
&\leq\frac{\varepsilon}{2}\sum_{m=1}^{M}a_m
+\sum_{m=M+1}^{\infty}a_m\\
&<\varepsilon.
\end{aligned}
\]
故 $F$ 右连续。

（3）仍选取上述 $M$。由于有限个分布函数 $F_1,\ldots,F_M$ 在负无穷处均趋于 $0$，当 $t$ 足够小时，
\[
0\leq F(t)
\leq\frac{\varepsilon}{2}\sum_{m=1}^{M}a_m
+\sum_{m=M+1}^{\infty}a_m
<\varepsilon.
\]
故 $\lim_{t\to-\infty}F(t)=0$。同理，当 $t$ 足够大时，$1-F_m(t)<\varepsilon/2$ 对所有 $1\leq m\leq M$ 成立，于是
\[
0\leq1-F(t)
=\sum_{m=1}^{\infty}a_m\bigl(1-F_m(t)\bigr)
<\varepsilon.
\]
故 $\lim_{t\to+\infty}F(t)=1$。

综上，$F$ 满足分布函数的三个本质条件，因此 $F$ 是一个分布函数。
\end{example}

\section{离散型随机变量}
\label{sec:discrete-random-variables}

本节介绍离散型随机变量。随机变量还包括连续型以及其他类型，不能将所有随机变量都归入离散型或连续型这两类。

\subsection*{离散型随机变量的定义和基本性质}

\begin{mydefinition}{离散型随机变量与分布列}{discrete-random-variable-pmf}
若随机变量 $X$ 的可能取值为有限个 $x_1,\ldots,x_s$ 或可列个 $x_1,x_2,\ldots$，则称 $X$ 为\textbf{离散型随机变量}。这些取值互不相同。令
\[
p(x_i):=P(X=x_i),
\]
称各取值及其对应概率构成的表为 $X$ 的\textbf{概率分布列}，简称\textbf{分布列}。函数 $p$ 也称为概率质量函数。

分布列可写成表格
\[
\begin{array}{c|ccccc}
X&x_1&x_2&\cdots&x_s&\cdots\\ \hline
P&p(x_1)&p(x_2)&\cdots&p(x_s)&\cdots
\end{array}
\]
或两行矩阵的形式
\[
X\sim
\begin{pmatrix}
x_1&x_2&\cdots&x_s&\cdots\\
p(x_1)&p(x_2)&\cdots&p(x_s)&\cdots
\end{pmatrix}.
\]
有限取值的情形去掉末尾的省略号。
\end{mydefinition}

\textbf{注：}\quad 有限取值或能够按递增次序列举的取值，通常从小到大排列。排列顺序不影响分布列，但各个不同取值应分别对应自己的概率。

\begin{proposition}{分布列的基本性质}{pmf-basic-properties}
离散型随机变量的分布列满足
\[
p(x_i)\geq0,\qquad \sum_i p(x_i)=1.
\]
其中求和遍历所有不同取值；有限情形为有限和，可列情形为无穷级数。
\end{proposition}

\begin{proof}
非负性由概率的非负性直接得到。各事件 $\{X=x_i\}$ 两两互斥，且其并集为 $\Omega$，故由可列可加性和规范性，
\[
\sum_i p(x_i)=\sum_iP(X=x_i)
=P\left(\bigcup_i\{X=x_i\}\right)=P(\Omega)=1.
\]
\end{proof}

\textbf{注：}\quad 反过来，给定互不相同的有限或可列个实数 $x_i$，以及满足 $q_i\geq0$、$\sum_iq_i=1$ 的数列，也能构造以 $q_i$ 为分布列的随机变量。具体可取
\[
\Omega=\{x_i\},\qquad
\mathscr F=\mathcal P(\Omega),\qquad
P(A)=\sum_{i:x_i\in A}q_i,\qquad
X(\omega)=\omega.
\]
这里 $\mathcal P(\Omega)$ 表示 $\Omega$ 的所有子集组成的集合族；于是 $P(X=x_i)=q_i$。
\newpage
\begin{proposition}{分布列与分布函数相互确定}{pmf-cdf-equivalence}
设 $X$ 为离散型随机变量，分布列为 $p(x_i)$，分布函数为 $F$。则
\[
F(t)=\sum_{i:x_i\leq t}p(x_i),\qquad t\in\mathbb R,
\]
并且
\[
p(x_i)=F(x_i)-F(x_i^-).
\]
\end{proposition}

\begin{proof}
事件 $\{X\leq t\}$ 是所有满足 $x_i\leq t$ 的互斥事件 $\{X=x_i\}$ 的并，故第一式由可列可加性得到。第二式由第 \ref{sec:random-variables-and-distribution-functions} 节的单点概率公式得到。
\end{proof}

\begin{example}[收集齐优惠券所需的张数]
有 $r$ 种不同优惠券，每次收集一张，每一种被收集到的概率均为 $1/r$，各次收集相互独立。设 $T$ 为首次收集齐全部 $r$ 种优惠券时累计收集的张数，求 $P(T=n)$。

当 $r=1$ 时，$P(T=1)=1$。以下设 $r\geq2$；显然 $n<r$ 时，$P(T=n)=0$。

先求 $P(T>m)$。设 $A_j^{(m)}$ 表示前 $m$ 次收集中没有第 $j$ 种优惠券，则
\[
\{T>m\}=\bigcup_{j=1}^rA_j^{(m)}.
\]
对任意不同的 $j_1,\ldots,j_k$，由独立性，
\[
P\left(A_{j_1}^{(m)}\cap\cdots\cap A_{j_k}^{(m)}\right)
=\left(1-\frac kr\right)^m.
\]
当 $m\geq r-1\geq1$ 时，容斥公式给出
\[
\begin{aligned}
P(T>m)
&=\sum_{k=1}^r(-1)^{k+1}\binom rk
\left(1-\frac kr\right)^m\\
&=\sum_{k=1}^{r-1}(-1)^{k+1}\binom rk
\left(1-\frac kr\right)^m,
\end{aligned}
\]
其中 $k=r$ 对应的项为 $0$。于是对 $n\geq r$，
\[
\begin{aligned}
P(T=n)
&=P(T>n-1)-P(T>n)\\
&=\sum_{k=1}^{r-1}(-1)^{k+1}\binom rk
\left[
\left(1-\frac kr\right)^{n-1}
-\left(1-\frac kr\right)^n
\right].
\end{aligned}
\]
\end{example}
\newpage
\subsection*{二项分布}

\begin{mydefinition}{二项分布}{binomial-distribution}
设 $n$ 为正整数，$0<p<1$。若随机变量 $X$ 取值于 $0,1,\ldots,n$，且
\[
P(X=k)=\binom nk p^k(1-p)^{n-k},
\qquad k=0,1,\ldots,n,
\]
则称 $X$ 服从参数为 $(n,p)$ 的\textbf{二项分布}，记为 $X\sim B(n,p)$。
\end{mydefinition}

\textbf{注：}\quad 二项分布也允许退化参数 $p=0$ 或 $p=1$：此时分别有 $P(X=0)=1$ 或 $P(X=n)=1$。下面涉及相邻概率比值的讨论均假设 $0<p<1$。

\textbf{注：}\quad 上式各项非负，且由二项式定理，
\[
\sum_{k=0}^n\binom nk p^k(1-p)^{n-k}
=\bigl(p+(1-p)\bigr)^n=1,
\]
所以它确实给出了一个分布列。

\textbf{注：}\quad $n=1$ 时，$B(1,p)$ 称为\textbf{伯努利分布}或\textbf{0--1 分布}，其分布列为
\[
P(X=1)=p,\qquad P(X=0)=1-p.
\]
一次只有“成功”和“失败”两种结果的试验，可用 $X=1$ 表示成功、$X=0$ 表示失败。

一般地，进行 $n$ 次相互独立的试验，每次成功的概率均为 $p$，则成功总次数 $X$ 服从 $B(n,p)$。由第一章的伯努利试验计算，恰好成功 $k$ 次有 $\binom nk$ 种成功位置，每一种位置安排的概率均为 $p^k(1-p)^{n-k}$。

\begin{proposition}{二项分布的单调性与最可能取值}{binomial-modes}
设 $X\sim B(n,p)$，$0<p<1$，令 $a=(n+1)p$。随着 $k$ 从 $0$ 增加到 $n$，概率 $P(X=k)$ 先不减后不增，最大值的位置为：
\begin{enumerate}
    \item 若 $a$ 不是整数，则唯一的最可能取值为 $\lfloor a\rfloor$；
    \item 若 $a=m$ 是整数，则 $m-1$ 和 $m$ 是两个并列的最可能取值，且 $1\leq m\leq n$。
\end{enumerate}
其中 $\lfloor a\rfloor$ 表示不大于 $a$ 的最大整数。
\end{proposition}

\begin{proof}
令 $q=1-p>0$。对 $k=1,\ldots,n$，相邻两项的比值为
\[
\frac{P(X=k)}{P(X=k-1)}
=\frac{\binom nk p^kq^{n-k}}
{\binom n{k-1}p^{k-1}q^{n-k+1}}
=\frac{(n-k+1)p}{kq}
=1+\frac{(n+1)p-k}{kq}.
\]
因为 $kq>0$，所以当 $k<a$、$k=a$、$k>a$ 时，该比值分别大于、等于、小于 $1$。

若 $a$ 非整数，相邻概率严格增加到 $\lfloor a\rfloor$ 处，此后严格减小；当 $\lfloor a\rfloor=0$ 或 $n$ 时，对应的一侧没有其他项。若 $a=m$ 为整数，则
\[
P(X=m-1)=P(X=m),
\]
这两项之前的概率严格增加，之后严格减小，因此两项同时达到最大值。
\end{proof}

\textbf{注：}\quad 最可能取值也称为最可能成功次数，相应的最大概率称为二项分布的\textbf{中心项}。$\lfloor(n+1)p\rfloor$ 总是一个最可能取值；当 $(n+1)p$ 为整数时，还有一个并列取值 $(n+1)p-1$。

\begin{example}[均匀硬币的中心项]
独立地抛一枚均匀硬币 $2n$ 次，设 $X_n$ 为正面出现的次数，则
\[
X_n\sim B\left(2n,\frac12\right).
\]
由于 $(2n+1)/2=n+1/2$ 不是整数，唯一的最可能正面次数为 $n$，且
\[
P(X_n=n)=\binom{2n}{n}\frac1{2^{2n}}
=\frac{(2n)!}{(n!)^2}\frac1{2^{2n}}.
\]
由斯特林（Stirling）公式
\[
n!\sim\sqrt{2\pi n}\left(\frac ne\right)^n,
\qquad n\to\infty,
\]
可得
\[
\begin{aligned}
P(X_n=n)
&\sim
\frac{\sqrt{4\pi n}(2n/e)^{2n}}
{2\pi n(n/e)^{2n}}\frac1{2^{2n}}\\
&=\frac1{\sqrt{\pi n}}\longrightarrow0.
\end{aligned}
\]
这里 $u_n\sim v_n$ 表示 $u_n/v_n\to1$。对所有 $0\leq k\leq2n$，
\[
0\leq P(X_n=k)\leq P(X_n=n).
\]
因此，对每个固定的非负整数 $k$，由夹逼定理，
\[
\lim_{n\to\infty}P(X_n=k)=0.
\]
\end{example}

\subsection*{泊松分布}

\begin{mydefinition}{泊松分布}{poisson-distribution}
设 $\lambda>0$。若随机变量 $X$ 取值于非负整数，且
\[
P(X=k)=\frac{\lambda^k}{k!}e^{-\lambda},
\qquad k=0,1,2,\ldots,
\]
则称 $X$ 服从参数为 $\lambda$ 的\textbf{泊松分布}，记为 $X\sim\mathcal P(\lambda)$。
\end{mydefinition}

\textbf{注：}\quad 各项概率非负。由指数函数的幂级数展开，
\[
\sum_{k=0}^{\infty}P(X=k)
=e^{-\lambda}\sum_{k=0}^{\infty}\frac{\lambda^k}{k!}
=e^{-\lambda}e^\lambda=1,
\]
所以该公式确实给出了一个概率分布列。

泊松随机变量常用来描述某段时间内或某个区域内事件发生的次数、对象的个数。

\begin{example}[一页上至少有一处印刷错误]
设一本书某一页上的印刷错误数 $X$ 服从参数为 $\lambda=1/2$ 的泊松分布。则
\[
P(X\geq1)=1-P(X=0)=1-e^{-1/2}\approx0.393469.
\]
\end{example}

\begin{example}[泊松粒子数的独立捕捉]
设某仪器在一定时间内发射的粒子总数 $\xi\sim\mathcal P(\lambda)$。给定 $\xi=n$ 时，每个粒子以相同概率 $p$ 被捕捉记录，且各粒子的捕捉结果相互独立，其中 $0<p<1$。记捕捉到的粒子总数为 $\eta$，证明 $\eta\sim\mathcal P(\lambda p)$。

令 $q=1-p$。给定 $\xi=n$ 时，$\eta$ 的条件分布为 $B(n,p)$。对任意非负整数 $k$，由全概率公式，
\[
\begin{aligned}
P(\eta=k)
&=\sum_{n=k}^{\infty}P(\xi=n)P(\eta=k\mid\xi=n)\\
&=\sum_{n=k}^{\infty}\frac{\lambda^n}{n!}e^{-\lambda}
\binom nk p^kq^{n-k}\\
&=\frac{e^{-\lambda}(\lambda p)^k}{k!}
\sum_{n=k}^{\infty}\frac{(\lambda q)^{n-k}}{(n-k)!}\\
&=\frac{e^{-\lambda}(\lambda p)^k}{k!}
\sum_{j=0}^{\infty}\frac{(\lambda q)^j}{j!}\\
&=\frac{(\lambda p)^k}{k!}e^{-\lambda}e^{\lambda q}
=\frac{(\lambda p)^k}{k!}e^{-\lambda p}.
\end{aligned}
\]
其中换元 $j=n-k$，所以 $\eta\sim\mathcal P(\lambda p)$。
\end{example}

\begin{proposition}{泊松分布的单调性与最可能取值}{poisson-modes}
设 $X\sim\mathcal P(\lambda)$。随着 $k$ 从 $0$ 增大，$P(X=k)$ 先不减后不增，并趋于 $0$。
\begin{enumerate}
    \item 若 $\lambda$ 不是整数，则唯一的最可能取值为 $\lfloor\lambda\rfloor$；
    \item 若 $\lambda=m$ 为正整数，则 $m-1$ 和 $m$ 是两个并列的最可能取值。
\end{enumerate}
\end{proposition}

\begin{proof}
对整数 $k\geq1$，
\[
\frac{P(X=k)}{P(X=k-1)}
=\frac{\lambda^ke^{-\lambda}/k!}
{\lambda^{k-1}e^{-\lambda}/(k-1)!}
=\frac{\lambda}{k}.
\]
所以当 $k<\lambda$、$k=\lambda$、$k>\lambda$ 时，相邻概率的比值分别大于、等于、小于 $1$，得到上述最可能取值。当 $0<\lambda<1$ 时，概率从 $k=0$ 开始严格减小，最大值在 $0$ 处。

再取正整数 $M>2\lambda$。对 $k\geq M$，
\[
P(X=k)=P(X=M-1)\prod_{j=M}^{k}\frac{\lambda}{j}
\leq P(X=M-1)\left(\frac12\right)^{k-M+1}
\longrightarrow0.
\]
\end{proof}

\begin{mytheorem}{泊松定理}{poisson-limit-theorem}
设 $X_n\sim B(n,p_n)$，$0<p_n<1$，且
\[
np_n\longrightarrow\lambda>0.
\]
则对每个固定的非负整数 $k$，
\[
\lim_{n\to\infty}P(X_n=k)=\frac{\lambda^k}{k!}e^{-\lambda}.
\]
\end{mytheorem}

\begin{proof}
记 $\lambda_n=np_n$，则 $\lambda_n\to\lambda$，且 $p_n=\lambda_n/n\to0$。对固定 $k$，当 $n\geq k$ 时，
\[
\begin{aligned}
P(X_n=k)
&=\binom nk p_n^k(1-p_n)^{n-k}\\
&=\frac1{k!}
\prod_{j=0}^{k-1}\left(1-\frac jn\right)
\lambda_n^k\left(1-\frac{\lambda_n}{n}\right)^{n-k}.
\end{aligned}
\]
$k=0$ 时乘积按空乘积等于 $1$ 处理。有限乘积趋于 $1$，$\lambda_n^k\to\lambda^k$。又由
\[
\frac{\log(1-u)}{u}\longrightarrow-1,\qquad u\to0^+,
\]
有
\[
(n-k)\log(1-p_n)
=(n-k)p_n\frac{\log(1-p_n)}{p_n}
\longrightarrow-\lambda.
\]
因此 $(1-p_n)^{n-k}\to e^{-\lambda}$，代回便得结论。
\end{proof}

\textbf{注：}\quad 当 $n$ 很大、$p$ 很小，且 $np$ 保持适当大小时，可用 $\mathcal P(np)$ 近似 $B(n,p)$：
\[
P(X=k)=\binom nk p^k(1-p)^{n-k}
\approx\frac{(np)^k}{k!}e^{-np}.
\]
$n>50$、$p<0.1$ 是经验提示，不是保证任意问题中近似误差都很小的严格判据。印刷错误数、某天进入商场的人数等计数问题，在符合独立性及小概率假设的模型下可考虑这种近似。

\begin{example}[死亡人数不超过五人的泊松近似]
某患病人群中，每人因该病死亡的概率为 $0.2\%$。在 $1000$ 人的结果相互独立且具有相同概率的模型下，设死亡人数为 $X$，则
\[
X\sim B(1000,0.002).
\]
精确概率为
\[
P(X\leq5)
=\sum_{k=0}^{5}\binom{1000}{k}0.002^k0.998^{1000-k}.
\]
令 $Z\sim\mathcal P(2)$，因为 $np=1000\times0.002=2$，所以
\[
P(X\leq5)\approx P(Z\leq5)
=e^{-2}\sum_{k=0}^{5}\frac{2^k}{k!}
\approx0.983436\approx0.9834.
\]
这是对原二项分布概率的近似，不能将 $X$ 的实际分布直接改写为泊松分布。
\end{example}

\textbf{注：}\quad 使用 WPS 表格，若 $Z\sim\mathcal P(\lambda)$，则函数返回值为
\[
\begin{aligned}
\texttt{POISSON}(k,\lambda,\texttt{TRUE})&=P(Z\leq k),\\
\texttt{POISSON}(k,\lambda,\texttt{FALSE})&=P(Z=k).
\end{aligned}
\]
例如，上例的近似概率可由 \texttt{POISSON(5, 2, TRUE)} 计算。

\subsection*{几何分布}

\begin{mydefinition}{几何分布}{geometric-distribution}
设 $0<p<1$。若随机变量 $X$ 取值于正整数，且
\[
P(X=k)=(1-p)^{k-1}p,\qquad k=1,2,\ldots,
\]
则称 $X$ 服从参数为 $p$ 的\textbf{几何分布}，记为 $X\sim G(p)$。
\end{mydefinition}

\textbf{注：}\quad 令 $q=1-p$。各项概率非负，且由等比级数求和，
\[
\sum_{k=1}^{\infty}q^{k-1}p=\frac{p}{1-q}=1,
\]
所以该公式确实是一个分布列。

几何分布对应独立重复伯努利试验中\textbf{首次成功所需的试验次数}，包括取得成功的那一次。这与第一章的首次成功概率计算一致。

\textbf{注：}\quad 对 $X\sim G(p)$ 和正整数 $k$，
\[
P(X\geq k)=\sum_{m=k}^{\infty}q^{m-1}p=q^{k-1},
\qquad
P(X>k)=q^k.
\]
也可以从试验意义理解：$X\geq k$ 表示前 $k-1$ 次全部失败，$X>k$ 表示前 $k$ 次全部失败。

\begin{mytheorem}{几何分布的无记忆性}{geometric-memorylessness}
设 $X$ 是取正整数值的随机变量，且对每个非负整数 $k$ 都有 $P(X>k)>0$。则 $X$ 服从某个参数 $0<p<1$ 的几何分布，当且仅当对任意正整数 $n$ 和非负整数 $k$，
\[
P(X=n+k\mid X>k)=P(X=n).
\]
\end{mytheorem}

\begin{proof}
先证必要性。设 $X\sim G(p)$，$q=1-p$，则
\[
P(X=n+k\mid X>k)
=\frac{P(X=n+k)}{P(X>k)}
=\frac{q^{n+k-1}p}{q^k}
=q^{n-1}p=P(X=n).
\]

再证充分性。令 $q=P(X>1)$。将无记忆性中的等式对 $n\geq2$ 求和，得到
\[
\begin{aligned}
P(X>k+1\mid X>k)
&=\sum_{n=2}^{\infty}P(X=k+n\mid X>k)\\
&=\sum_{n=2}^{\infty}P(X=n)=q.
\end{aligned}
\]
所以
\[
P(X>k+1)=qP(X>k),\qquad k\geq0.
\]
由 $P(X>0)=1$，递推得到 $P(X>k)=q^k$。前提保证 $q>0$；又由于 $X$ 取有限的正整数值，
\[
P(X>k)\longrightarrow0,
\]
故 $q<1$。因此令 $p=1-q\in(0,1)$，对每个正整数 $n$，
\[
P(X=n)=P(X>n-1)-P(X>n)
=q^{n-1}-q^n=q^{n-1}p,
\]
即 $X\sim G(p)$。
\end{proof}

\textbf{注：}\quad 无记忆性表示：已知前 $k$ 次试验都失败后，从此刻起再等待首次成功的次数，与从头开始的等待次数具有相同分布。这里的 $P(X>k)>0$ 保证条件概率有定义；无记忆性说的是概率规律，不表示试验者具有或缺乏学习能力。

\subsection*{帕斯卡分布}

\begin{mydefinition}{帕斯卡分布}{pascal-distribution}
设 $n$ 为正整数，$0<p<1$。若随机变量 $X$ 取值于 $n,n+1,\ldots$，且
\[
P(X=k)=\binom{k-1}{n-1}p^n(1-p)^{k-n},
\qquad k=n,n+1,\ldots,
\]
则称 $X$ 服从参数为 $(n,p)$ 的\textbf{帕斯卡分布}，记为 $X\sim\mathrm{Pa}(n,p)$。
\end{mydefinition}

\textbf{注：}\quad 验证归一化时，可用 $|x|<1$ 范围内的泰勒展开：
\[
\begin{aligned}
(1-x)^{-n}
&=\sum_{m=0}^{\infty}\binom{-n}{m}(-x)^m\\
&=\sum_{m=0}^{\infty}
\frac{n(n+1)\cdots(n+m-1)}{m!}x^m\\
&=\sum_{m=0}^{\infty}\binom{n+m-1}{m}x^m\\
&=\sum_{k=n}^{\infty}\binom{k-1}{n-1}x^{k-n}.
\end{aligned}
\]
这里 $\binom{-n}{m}=(-1)^m\binom{n+m-1}{m}$，$m=0$ 时乘积等于 $1$。令 $x=1-p$，得到
\[
\sum_{k=n}^{\infty}\binom{k-1}{n-1}p^n(1-p)^{k-n}
=p^n\bigl[1-(1-p)\bigr]^{-n}=1.
\]
结合非负性，说明上述公式确实给出了概率分布列。

\textbf{注：}\quad 帕斯卡分布对应独立重复伯努利试验中\textbf{取得第 $n$ 次成功时的总试验次数}。第 $k$ 次必须成功，而前 $k-1$ 次恰有 $n-1$ 次成功，因此其概率为上述分布列。特别地，$\mathrm{Pa}(1,p)=G(p)$。

\begin{mydefinition}{负二项分布}{negative-binomial-distribution}
若 $X\sim\mathrm{Pa}(n,p)$，令 $Y=X-n$，则 $Y$ 表示取得第 $n$ 次成功时已经发生的\textbf{失败次数}，取值于非负整数。其分布列为
\[
\begin{aligned}
P(Y=k)&=P(X=n+k)\\
&=\binom{n+k-1}{n-1}p^n(1-p)^k\\
&=\binom{n+k-1}{k}p^n(1-p)^k,\qquad k=0,1,\ldots.
\end{aligned}
\]
称 $Y$ 服从参数为 $(n,p)$ 的\textbf{负二项分布}，记为 $Y\sim\mathrm{NB}(n,p)$。
\end{mydefinition}

\textbf{注：}\quad 广义二项式系数满足
\[
\binom{-n}{k}
=\frac{(-n)(-n-1)\cdots(-n-k+1)}{k!}
=(-1)^k\binom{n+k-1}{k},
\]
所以也可写为
\[
P(Y=k)=\binom{-n}{k}p^n(p-1)^k.
\]
这里区分两种计数方式：帕斯卡分布记录总试验次数，负二项分布记录失败次数；两者相差固定的 $n$ 次成功。

\begin{example}[巴拿赫火柴盒问题]
一个人左右口袋各放一盒火柴，每盒最初有 $N$ 根，$N$ 为正整数。每次需要火柴时，相互独立地以概率 $1/2$ 选择一盒；若盒内有火柴，就取出一根。求他首次选到一个空盒、发现该盒已用光时，另一盒恰好剩下 $k$ 根的概率，其中 $0\leq k\leq N$。

设 $A_k$ 为首次发现右盒为空时，左盒恰好剩下 $k$ 根的事件。由于左右对称，两种首次发现空盒的事件互斥，所求概率为 $2P(A_k)$。

$A_k$ 发生时，这是第 $N+1$ 次选择右盒，而此前共选择左盒 $N-k$ 次。将选择右盒视为成功，则
\[
T\sim\mathrm{Pa}\left(N+1,\frac12\right),
\]
其中 $T$ 表示第 $N+1$ 次选择右盒的时刻。$A_k$ 等价于 $T=2N-k+1$；此前选择左盒的次数不超过 $N$，因而还没有发现左盒为空。于是
\[
\begin{aligned}
P(A_k)
&=\binom{2N-k}{N}
\left(\frac12\right)^{N+1}
\left(\frac12\right)^{N-k}\\
&=\binom{2N-k}{N}\frac1{2^{2N-k+1}}.
\end{aligned}
\]
因此所求概率为
\[
\boxed{\,P(\text{另一盒剩下 }k\text{ 根})
=\binom{2N-k}{N}\frac1{2^{2N-k}},\quad 0\leq k\leq N.\,}
\]

\textbf{注：}\quad 停止时刻是“首次选到空盒并发现已经用光”，不是“刚取走某盒的最后一根火柴”；区分这两个时刻才能得到正确的分布。
\end{example}

\section{连续型随机变量}
\label{sec:continuous-random-variables}

\subsection*{从小区间频率到概率密度}

离散型随机变量的区间概率可以按取值求和。对于寿命等通常以区间为取值范围的量，可以考虑区间概率。

将区间划分为 $a=x_0<x_1<\cdots<x_n=b$，记 $\Delta x_i=x_i-x_{i-1}$。由互斥事件的可加性，
\[
P(a<X\leq b)=\sum_{i=1}^{n}P(x_{i-1}<X\leq x_i).
\]
做 $M$ 次独立同分布试验，若有 $m_i$ 次取值落在 $(x_{i-1},x_i]$ 中，则其频率为 $m_i/M$，每单位区间长度的频率为 $m_i/(M\Delta x_i)$。用函数 $f$ 描述这种密度时，直观上有
\[
\begin{aligned}
P(a<X\leq b)
&\approx\sum_{i=1}^{n}\frac{m_i}{M}\\
&=\sum_{i=1}^{n}\frac{m_i}{M\Delta x_i}\Delta x_i\\
&\approx\sum_{i=1}^{n}f(x_i)\Delta x_i
\approx\int_a^b f(x)\,dx.
\end{aligned}
\]
这解释了为什么区间概率对应密度曲线下的面积，但不是密度存在性的证明。

\textbf{注：}\quad 取值不可数，或分布函数连续，都不能单独保证随机变量具有概率密度；本节的定义以存在下面的积分表示为准。

\subsection*{连续型随机变量与概率密度}

\begin{mydefinition}{连续型随机变量与概率密度}{continuous-random-variable-density}
设 $X$ 为随机变量。若存在非负可测函数 $f:\mathbb R\to[0,\infty)$，使得对任意 $B\in\mathscr B(\mathbb R)$，
\[
P(X\in B)=\int_B f(x)\,dx,
\]
则称 $X$ 为\textbf{连续型随机变量}，称 $f$ 为 $X$ 的\textbf{概率密度函数}，简称\textbf{密度函数}或\textbf{密度}。
\end{mydefinition}

\textbf{注：}\quad 定义中一般集合上的积分是勒贝格积分。本节常见分段函数的区间概率先按熟悉的定积分或反常积分计算。直观上，$f(x)\Delta x$ 近似一个小区间的概率；$f(x)$ 本身不是 $P(X=x)$，也不要求 $f(x)\leq1$。

\begin{proposition}{概率密度的基本性质}{density-basic-properties}
若 $f$ 是连续型随机变量 $X$ 的密度，则
\[
f(x)\geq0,\qquad
\int_{\mathbb R}f(x)\,dx
=\int_{-\infty}^{+\infty}f(x)\,dx=1.
\]
\end{proposition}

\begin{proof}
非负性由定义得到。取 $B=\mathbb R$，则
\[
\int_{\mathbb R}f(x)\,dx=P(X\in\mathbb R)=P(\Omega)=1.
\]
\end{proof}

\textbf{注：}\quad 反过来，非负可测函数 $g$ 若满足 $\int_{\mathbb R}g(x)\,dx=1$，就能成为某个随机变量的密度。可以取
\[
\Omega=\mathbb R,\qquad
\mathscr F=\mathscr B(\mathbb R),\qquad
P(B)=\int_Bg(x)\,dx,\qquad
Y(\omega)=\omega.
\]

\begin{example}[分段密度的参数、区间概率与分布函数]
设 $X$ 的概率密度为
\[
f(x)=
\begin{cases}
\dfrac18,&0<x<2,\\[0.4em]
kx,&2\leq x<4,\\
0,&\text{其他}.
\end{cases}
\]
求 $k$、$P(1<X\leq3)$ 以及分布函数 $F$。

（1）由密度的规范性，
\[
1=\int_0^2\frac18\,dx+\int_2^4kx\,dx
=\frac14+6k,
\]
故 $k=1/8$，也满足非负性要求。

（2）区间跨越密度的分段点 $2$，因此
\[
P(1<X\leq3)
=\int_1^2\frac18\,dx+\int_2^3\frac{x}{8}\,dx
=\frac18+\frac5{16}=\frac7{16}.
\]

（3）由 $F(t)=\int_{-\infty}^t f(x)\,dx$ 分段计算。特别地，$2\leq t<4$ 时，
\[
F(t)=\int_0^2\frac18\,dx+\int_2^t\frac{x}{8}\,dx
=\frac14+\frac{t^2-4}{16}=\frac{t^2}{16}.
\]
综上，
\[
F(t)=
\begin{cases}
0,&t<0,\\[0.3em]
\dfrac t8,&0\leq t<2,\\[0.5em]
\dfrac{t^2}{16},&2\leq t<4,\\[0.5em]
1,&t\geq4.
\end{cases}
\]
在 $0,2,4$ 处左右极限与函数值均相同，故 $F$ 连续。
\end{example}

\subsection*{密度与分布函数的关系}

\begin{proposition}{连续型分布的基本性质}{continuous-distribution-properties}
设 $X$ 为连续型随机变量，密度与分布函数分别为 $f$、$F$。则：
\begin{enumerate}
    \item 对任意 $t\in\mathbb R$，
    \[
    F(t)=\int_{-\infty}^t f(x)\,dx;
    \]
    \item $F$ 连续，且几乎处处可微，并且
    \[
    F'(t)=f(t)\qquad\text{几乎处处成立};
    \]
    在 $f$ 连续的点 $t$，由积分上限函数求导可得 $F'(t)=f(t)$；
    \item 若 Borel 集 $B$ 的勒贝格测度 $m(B)=0$，则
    \[
    P(X\in B)=0.
    \]
    特别地，对任意实数 $c$，$P(X=c)=0$；
    \item 对给定的随机变量 $X$，分布函数唯一，密度则只在相差一个零测集的意义下唯一。即若 $f_1,f_2$ 都是其密度，则
    \[
    m\bigl(\{x:f_1(x)\ne f_2(x)\}\bigr)=0.
    \]
\end{enumerate}
\end{proposition}

\begin{proof}
（1）由定义，
\[
F(t)=P(X\leq t)=P(X\in(-\infty,t])
=\int_{-\infty}^t f(x)\,dx.
\]

（2）在密度连续的点，利用微积分基本定理对积分上限求导。

（3）利用勒贝格积分在零测集上的积分为 $0$，得到
\[
P(X\in B)=\int_Bf(x)\,dx=0.
\]
单点集合的测度为 $0$，所以 $P(X=c)=0$。

（4）分布函数由 $F(t)=P(X\leq t)$ 唯一确定。若 $f_1,f_2$ 都是 $X$ 的密度，则对任意 Borel 集 $B$，
\[
\int_Bf_1(x)\,dx=\int_Bf_2(x)\,dx.
\]
由积分唯一性，两个密度几乎处处相等。
\end{proof}

\textbf{注：}\quad “几乎处处”是指除去一个勒贝格测度为 $0$ 的集合后成立；这里的测度是区间长度概念的推广。改变密度在单个点或零测集上的非负取值，不改变其所描述的分布。求导等式须按上述条件和几乎处处的意义理解。

\textbf{注：}\quad 因为各端点的单点概率为 $0$，对 $a<b$，
\[
\begin{aligned}
P(a<X\leq b)&=P(a\leq X\leq b)\\
&=P(a<X<b)=P(a\leq X<b)\\
&=\int_a^b f(x)\,dx=F(b)-F(a).
\end{aligned}
\]
本节中开闭端点不影响区间概率；离散型或一般随机变量的计算仍须检查端点质量。

\subsection*{利用密度与分布函数计算概率}

\begin{example}[五个电子管中恰有两个需要更换]
某类电子管的寿命 $X$（单位：小时）具有密度
\[
f(x)=
\begin{cases}
\dfrac{100}{x^2},&x>100,\\[0.4em]
0,&x\leq100.
\end{cases}
\]
现有 $5$ 个具有相同寿命分布且相互独立的电子管，求 $150$ 小时内恰有 $2$ 个原装电子管因寿命耗尽需要更换的概率。

单个电子管在 $150$ 小时内需要更换的概率为
\[
p=P(X\leq150)
=\int_{100}^{150}\frac{100}{x^2}\,dx
=\left[-\frac{100}{x}\right]_{100}^{150}
=\frac13.
\]
记需要更换的原装电子管个数为 $Y$。由寿命的独立性，
\[
Y\sim B\left(5,\frac13\right),
\]
所以
\[
P(Y=2)=\binom52\left(\frac13\right)^2
\left(\frac23\right)^3=\frac{80}{243}.
\]
\end{example}

\begin{example}[由分布函数求密度]
设连续型随机变量 $X$ 的分布函数为
\[
F(t)=
\begin{cases}
0,&t<0,\\
t^2,&0\leq t<1,\\
1,&t\geq1.
\end{cases}
\]
在各个开区间上求导，可选择一个密度版本为
\[
f(x)=
\begin{cases}
2x,&0<x<1,\\
0,&\text{其他}.
\end{cases}
\]
因为 $F$ 在 $1$ 处不可微，$f(1)=0$ 是选定的密度值，不是该点的导数；零测集上的取值可以另作非负指定。并且
\[
\int_0^1 2x\,dx=1,
\]
将此密度从 $-\infty$ 积分到 $t$，恰好还原上面的 $F(t)$。

\textbf{注：}\quad 可简写为“$f(x)=2x,\ 0<x<1$，其余处为 $0$”；不能把区间外的零值条件省略成对所有实数都有 $f(x)=2x$。
\end{example}

\subsection*{取值区间与密度的零值部分}

\textbf{注：}\quad 若 $P(a<X\leq b)=1$，称 $(a,b]$ 为 $X$ 的取值区间。这里表示 $X$ 以概率 $1$ 落在此区间，不必是每个样本点上的严格取值范围，也不保证区间内每一处密度都为正。

由 $f\geq0$ 和区间外的积分为 $0$，密度在区间外几乎处处为 $0$，因此可以选取一个密度版本，使
\[
f(x)=0,\qquad x\notin(a,b].
\]
类似地，若 $P(X\geq0)=1$，则称 $X$ 取值非负，可以选取密度满足
\[
f(x)=0,\qquad x<0.
\]

\subsection*{均匀分布}

\begin{mydefinition}{均匀分布}{uniform-distribution}
设 $a<b$。若随机变量 $X$ 的密度可以取为
\[
f(x)=
\begin{cases}
\dfrac1{b-a},&a<x<b,\\[0.4em]
0,&\text{其他},
\end{cases}
\]
则称 $X$ 服从区间 $(a,b)$ 上的\textbf{均匀分布}，记为 $X\sim U(a,b)$。
\end{mydefinition}

\textbf{注：}\quad 此函数非负，且
\[
\int_{\mathbb R}f(x)\,dx
=\int_a^b\frac1{b-a}\,dx
=\frac{b-a}{b-a}=1,
\]
所以确实是密度。对包含于 $(a,b)$ 中、长度为 $\ell$ 的子区间 $(c,c+\ell)$，
\[
P(c<X<c+\ell)
=\int_c^{c+\ell}\frac1{b-a}\,dx
=\frac{\ell}{b-a}.
\]
概率只取决于子区间的长度，与位置无关。这就是“均匀”的含义，也与区间上的几何概型一致。

\begin{example}[等公交车不超过五分钟]
公交车从 $7$ 点起每隔 $15$ 分钟到站一次。乘客到站时间服从 $7$ 点至 $7$ 点 $30$ 分之间的均匀分布，求等车时间不超过 $5$ 分钟的概率。

设 $X$ 为乘客到站时距 $7$ 点的分钟数，则 $X\sim U(0,30)$。在本时间段内公交车于 $0,15,30$ 分钟到站，所求事件为
\[
A=\{X=0\}\cup\{10\leq X\leq15\}\cup\{25\leq X\leq30\}.
\]
上述事件互斥，且 $P(X=0)=0$，因此
\[
\begin{aligned}
P(A)
&=P(X=0)+P(10\leq X\leq15)+P(25\leq X\leq30)\\
&=0+\int_{10}^{15}\frac1{30}\,dx+\int_{25}^{30}\frac1{30}\,dx\\
&=\frac5{30}+\frac5{30}=\frac13.
\end{aligned}
\]
\end{example}

\subsection*{指数分布}

\begin{mydefinition}{指数分布}{exponential-distribution}
设 $\lambda>0$。若随机变量 $X$ 的密度可以取为
\[
f(x)=
\begin{cases}
\lambda e^{-\lambda x},&x\geq0,\\
0,&x<0,
\end{cases}
\]
则称 $X$ 服从参数为 $\lambda$ 的\textbf{指数分布}，记为 $X\sim\operatorname{Exp}(\lambda)$。这里 $\lambda$ 使用率参数的约定。
\end{mydefinition}

\textbf{注：}\quad 上述函数非负，且
\[
\begin{aligned}
\int_{\mathbb R}f(x)\,dx
&=\lim_{R\to\infty}\int_0^R\lambda e^{-\lambda x}\,dx\\
&=\lim_{R\to\infty}\bigl(1-e^{-\lambda R}\bigr)=1,
\end{aligned}
\]
所以确实是密度。

\begin{proposition}{指数分布的分布函数与尾概率}{exponential-cdf}
对连续型随机变量 $X$，$X\sim\operatorname{Exp}(\lambda)$ 当且仅当
\[
F(t)=
\begin{cases}
1-e^{-\lambda t},&t\geq0,\\
0,&t<0.
\end{cases}
\]
等价地，其尾概率为
\[
P(X>t)=
\begin{cases}
e^{-\lambda t},&t\geq0,\\
1,&t<0.
\end{cases}
\]
\end{proposition}

\begin{proof}
先证必要性。若 $X$ 服从上述指数分布，则 $t<0$ 时 $F(t)=0$；$t\geq0$ 时，
\[
F(t)=\int_0^t\lambda e^{-\lambda x}\,dx
=1-e^{-\lambda t}.
\]
再由 $P(X>t)=1-F(t)$ 得到尾概率。

再证充分性。若 $F$ 具有上述形式，则将候选密度
\[
g(x)=\lambda e^{-\lambda x},\quad x\geq0,
\qquad g(x)=0,\quad x<0
\]
积分到任意上限 $t$，恰好得到 $F(t)$。故 $g$ 是 $X$ 的一个密度版本，$X\sim\operatorname{Exp}(\lambda)$。
\end{proof}

\begin{mytheorem}{指数分布的无记忆性}{exponential-memorylessness}
设 $X$ 为取值非负的连续型随机变量，且对每个 $s>0$，$P(X>s)>0$。则 $X$ 服从某个参数 $\lambda>0$ 的指数分布，当且仅当对任意 $s,t>0$，
\[
P(X>s+t\mid X>s)=P(X>t).
\]
这一性质也称为\textbf{无效性}。
\end{mytheorem}

\begin{proof}
\textbf{必要性。}\quad 设 $X\sim\operatorname{Exp}(\lambda)$。因为 $s,t>0$，
\[
\begin{aligned}
P(X>s+t\mid X>s)
&=\frac{P(X>s+t)}{P(X>s)}\\
&=\frac{e^{-\lambda(s+t)}}{e^{-\lambda s}}\\
&=e^{-\lambda t}=P(X>t).
\end{aligned}
\]

\textbf{充分性。}\quad 令 $G(t)=P(X>t)$，$t>0$。无记忆性给出
\[
\begin{aligned}
G(s+t)
&=P(X>s+t\mid X>s)P(X>s)\\
&=G(t)G(s).
\end{aligned}
\]
由于 $X$ 连续型，$F$ 连续，故 $G=1-F$ 连续；前提保证 $G(t)>0$。由后面的引理，存在 $b\in\mathbb R$，使
\[
G(t)=e^{bt},\qquad t>0.
\]
由 $G(t)\leq1$ 可知 $b\leq0$。若 $b=0$，则所有 $t>0$ 都有 $G(t)=1$，这与
\[
\lim_{t\to\infty}G(t)
=1-\lim_{t\to\infty}F(t)=0
\]
矛盾。因此 $b<0$。令 $\lambda=-b>0$，便有
\[
P(X>t)=e^{-\lambda t},\qquad t>0.
\]
又因 $X$ 非负且单点概率为 $0$，$P(X>0)=1$；故其完整尾概率具有上一命题中的形式，$X\sim\operatorname{Exp}(\lambda)$。
\end{proof}

\begin{mylemma}{连续乘法方程}{continuous-multiplicative-equation}
若连续函数 $G:(0,\infty)\to(0,\infty)$ 满足
\[
G(s+t)=G(s)G(t),\qquad s,t>0,
\]
则存在 $b\in\mathbb R$，使 $G(t)=e^{bt}$ 对所有 $t>0$ 成立。
\end{mylemma}

\begin{proof}
令 $h(t)=\log G(t)$。因为 $G$ 正且连续，$h$ 连续，并满足
\[
h(s+t)=h(s)+h(t).
\]

先对正整数倍使用归纳法。$n=1$ 时显然成立；若 $h(nx)=nh(x)$，则
\[
h((n+1)x)=h(nx)+h(x)=(n+1)h(x).
\]
故对任意正整数 $n$，
\[
h(nx)=nh(x),\qquad x>0.
\]

令 $x=1/n$，得到
\[
h(1)=nh(1/n),
\qquad h(1/n)=\frac{h(1)}n.
\]
于是对任意正有理数 $q=m/n$，
\[
\begin{aligned}
h(q)&=h(m/n)\\
&=mh(1/n)\\
&=\frac mn h(1)=qh(1).
\end{aligned}
\]

最后，对任意 $t>0$，取正有理数列 $q_j\to t$。由连续性，
\[
\begin{aligned}
h(t)&=\lim_{j\to\infty}h(q_j)\\
&=\lim_{j\to\infty}q_jh(1)=th(1).
\end{aligned}
\]
令 $b=h(1)$，则 $\log G(t)=bt$，所以 $G(t)=e^{bt}$。
\end{proof}

\textbf{注：}\quad 若 $X$ 表示寿命，无记忆性表示：在已正常工作 $s$ 小时的条件下，再存活超过 $t$ 小时的概率，与新仪器存活超过 $t$ 小时的概率相同。相同的是剩余寿命的分布，不是说已过去的工作时间从总寿命中消失。条件 $P(X>s)>0$ 保证条件概率有定义。

\begin{example}[邮局中最后办完事情的人]
邮局有两位职员。$A$ 到达时，$B$、$C$ 正分别接受服务；其中一人办完后，$A$ 立即开始接受服务。假定各人的服务时间相互独立，均服从同一 $\operatorname{Exp}(\lambda)$ 分布，求 $A$ 在三人中最后办完的概率。

由无记忆性，$A$ 开始服务时，另一位尚未办完者的剩余服务时间 $R$ 仍服从 $\operatorname{Exp}(\lambda)$。$A$ 的服务时间 $T$ 也服从同一分布，且与 $R$ 独立。此时已有一人办完，因此
\[
\{A\text{ 最后办完}\}=\{T>R\}.
\]
独立同分布给出对称性，
\[
P(T>R)=P(R>T).
\]
两者独立且为连续型，$P(T=R)=0$，所以
\[
P(T>R)+P(R>T)=1,
\qquad
P(A\text{ 最后办完})=\frac12.
\]

\textbf{注：}\quad 该模型要求服务时间相互独立；仅有相同的指数边缘分布，并不足以推出上述比较结论。
\end{example}

\begin{example}[第一位顾客的到达时刻]
设 $N(t)$ 表示商场从开门时刻 $0$ 到时刻 $t$ 的顾客到达人数，且对每个 $t>0$，
\[
N(t)\sim\mathcal P(\lambda t),\qquad \lambda>0.
\]
令 $S$ 为第一位顾客的到达时刻。对 $t>0$，
\[
\{S>t\}=\{N(t)=0\},
\]
所以
\[
\begin{aligned}
P(S>t)
&=P(N(t)=0)\\
&=e^{-\lambda t}.
\end{aligned}
\]
结合到达时刻非负可知，$S\sim\operatorname{Exp}(\lambda)$。
\end{example}

\subsection*{\texorpdfstring{$\Gamma$}{Gamma} 分布}

\begin{mydefinition}{Gamma 函数与 Gamma 分布}{gamma-function-and-distribution}
对 $\alpha>0$，定义
\[
\Gamma(\alpha)=\int_0^\infty x^{\alpha-1}e^{-x}\,dx,
\]
称为\textbf{Gamma 函数}或 $\Gamma$ 函数。

设 $\alpha>0$、$\lambda>0$。若随机变量 $X$ 的密度可以取为
\[
f(x)=
\begin{cases}
\dfrac{\lambda^\alpha}{\Gamma(\alpha)}x^{\alpha-1}e^{-\lambda x},
&x>0,\\[0.5em]
0,&x\leq0,
\end{cases}
\]
则称 $X$ 服从参数为 $(\alpha,\lambda)$ 的\textbf{Gamma 分布}，记为 $X\sim\Gamma(\alpha,\lambda)$。其中 $\alpha$ 为形状参数，$\lambda$ 为率参数；正半轴上的密度也可写成
\[
f(x)=\frac{(\lambda x)^{\alpha-1}}{\Gamma(\alpha)}
\lambda e^{-\lambda x},\qquad x>0.
\]
\end{mydefinition}

\textbf{注：}\quad 当 $\alpha$ 为正整数时，该分布称为\textbf{埃尔朗分布}。特别地，$\Gamma(1,\lambda)=\operatorname{Exp}(\lambda)$。

\textbf{补充说明：}\quad 定义中的反常积分对 $\alpha>0$ 收敛。在原点附近，
\[
0\leq\int_0^1x^{\alpha-1}e^{-x}\,dx
\leq\int_0^1x^{\alpha-1}\,dx=\frac1\alpha.
\]
在 $[1,\infty)$ 上，函数 $x^{\alpha-1}e^{-x/2}$ 连续且趋于 $0$，因而有界。存在常数 $C_\alpha>0$，使
\[
0\leq x^{\alpha-1}e^{-x}\leq C_\alpha e^{-x/2},
\qquad x\geq1.
\]
右侧在 $[1,\infty)$ 可积，故 $\Gamma(\alpha)$ 有限；又因被积函数为正，$\Gamma(\alpha)>0$。

\begin{proposition}{Gamma 函数的基本性质}{gamma-function-properties}
Gamma 函数满足：
\begin{enumerate}
    \item 对 $\alpha>0$，$\Gamma(\alpha+1)=\alpha\Gamma(\alpha)$；
    \item $\Gamma(1)=1$，$\Gamma(1/2)=\sqrt\pi$；
    \item 对正整数 $n$，$\Gamma(n)=(n-1)!$。
\end{enumerate}
\end{proposition}

\begin{proof}
\textbf{（1）递推公式。}\quad 使用分部积分。先在 $[\varepsilon,R]$ 上计算，其中 $0<\varepsilon<R$：
\[
\begin{aligned}
\int_\varepsilon^R x^\alpha e^{-x}\,dx
&=\left[-x^\alpha e^{-x}\right]_\varepsilon^R
+\alpha\int_\varepsilon^R x^{\alpha-1}e^{-x}\,dx.
\end{aligned}
\]
因为 $\alpha>0$，
\[
\lim_{\varepsilon\downarrow0}\varepsilon^\alpha e^{-\varepsilon}=0,
\qquad
\lim_{R\to\infty}R^\alpha e^{-R}=0.
\]
令 $\varepsilon\downarrow0$、$R\to\infty$，由两侧反常积分的收敛性，
\[
\Gamma(\alpha+1)=\alpha\Gamma(\alpha).
\]

\textbf{（2）两个特殊值。}\quad 首先，
\[
\Gamma(1)=\int_0^\infty e^{-x}\,dx=1.
\]
对于 $\Gamma(1/2)$，作换元 $x=y^2/2$。此时
\[
dx=y\,dy,\qquad x^{-1/2}=\frac{\sqrt2}{y},\qquad y>0.
\]
因此
\[
\begin{aligned}
\Gamma(1/2)
&=\int_0^\infty x^{-1/2}e^{-x}\,dx\\
&=\sqrt2\int_0^\infty e^{-y^2/2}\,dy\\
&=\frac{\sqrt2}{2}\int_{-\infty}^{\infty}e^{-y^2/2}\,dy\\
&=\frac{\sqrt2}{2}\sqrt{2\pi}=\sqrt\pi.
\end{aligned}
\]
最后一步使用的高斯积分推导如下。

\textbf{（3）正整数值。}\quad $n=1$ 时，$\Gamma(1)=1=0!$。若 $\Gamma(n)=(n-1)!$，由（1），
\[
\Gamma(n+1)=n\Gamma(n)=n(n-1)!=n!.
\]
故由归纳法，对全部正整数 $n$，$\Gamma(n)=(n-1)!$。
\end{proof}

\textbf{高斯积分：}\quad
$\int_{\mathbb R}e^{-y^2/2}\,dy=\sqrt{2\pi}$。令
\[
I_R=\int_{-R}^Re^{-y^2/2}\,dy,
\qquad
D_R=\{(u,v):u^2+v^2\leq R^2\}.
\]
在有界区域上作二重积分，并换为极坐标，
\[
\begin{aligned}
\iint_{D_R}e^{-(u^2+v^2)/2}\,du\,dv
&=\int_0^{2\pi}\int_0^R e^{-r^2/2}r\,dr\,d\theta\\
&=2\pi(1-e^{-R^2/2}).
\end{aligned}
\]
又因 $D_R\subseteq[-R,R]^2\subseteq D_{\sqrt2R}$，被积函数非负，且矩形上的积分可分离为乘积，得到
\[
2\pi(1-e^{-R^2/2})
\leq I_R^2
\leq2\pi(1-e^{-R^2}).
\]
令 $R\to\infty$，由夹逼定理 $I_R^2\to2\pi$。因为 $I_R>0$，
\[
\int_{-\infty}^{\infty}e^{-y^2/2}\,dy
=\lim_{R\to\infty}I_R=\sqrt{2\pi}.
\]

\begin{proposition}{Gamma 密度的归一化}{gamma-density-normalization}
上述 Gamma 分布定义中的 $f$ 确实是概率密度函数。
\end{proposition}

\begin{proof}
由 $\alpha,\lambda,\Gamma(\alpha)>0$，$f$ 非负。作换元 $u=\lambda x$，有
\[
\begin{aligned}
\int_{\mathbb R}f(x)\,dx
&=\frac1{\Gamma(\alpha)}
\int_0^\infty(\lambda x)^{\alpha-1}\lambda e^{-\lambda x}\,dx\\
&=\frac1{\Gamma(\alpha)}
\int_0^\infty u^{\alpha-1}e^{-u}\,du\\
&=\frac{\Gamma(\alpha)}{\Gamma(\alpha)}=1.
\end{aligned}
\]
\end{proof}

\begin{example}[第 $n$ 位顾客的到达时刻]
仍设商场在 $[0,t]$ 内的到达人数满足
\[
N(t)\sim\mathcal P(\lambda t),\qquad t>0.
\]
令 $S_n$ 为第 $n$ 位顾客的到达时刻，$n\geq1$。证明 $S_n\sim\Gamma(n,\lambda)$。

先求分布函数，再求导。对 $t>0$，
\[
\{S_n\leq t\}=\{N(t)\geq n\}.
\]
所以
\[
\begin{aligned}
F_{S_n}(t)
&=P(N(t)\geq n)\\
&=1-P(N(t)\leq n-1)\\
&=1-e^{-\lambda t}\sum_{k=0}^{n-1}\frac{(\lambda t)^k}{k!}.
\end{aligned}
\]
$t\leq0$ 时取 $F_{S_n}(t)=0$，与右侧在 $0$ 处的极限一致。

对 $t>0$ 求导，将乘积的两部分分别列出：
\[
\begin{aligned}
f_{S_n}(t)
&=e^{-\lambda t}\sum_{k=0}^{n-1}
\frac{\lambda^{k+1}t^k}{k!}
-e^{-\lambda t}\sum_{k=1}^{n-1}
\frac{\lambda^kt^{k-1}}{(k-1)!}\\
&=e^{-\lambda t}\left[
\sum_{k=0}^{n-1}\frac{\lambda^{k+1}t^k}{k!}
-\sum_{j=0}^{n-2}\frac{\lambda^{j+1}t^j}{j!}
\right]\\
&=\frac{\lambda^nt^{n-1}}{(n-1)!}e^{-\lambda t}\\
&=\frac{(\lambda t)^{n-1}}{\Gamma(n)}
\lambda e^{-\lambda t}.
\end{aligned}
\]
前两式中两个求和的共有项逐项抵消；$n=1$ 时第二个求和为空和，按 $0$ 处理。取 $t\leq0$ 时密度为 $0$，便得到 Gamma 密度，故
\[
S_n\sim\Gamma(n,\lambda).
\]
特别地，$S_1\sim\Gamma(1,\lambda)=\operatorname{Exp}(\lambda)$，与第一位顾客的结果一致。
\end{example}

\subsection*{贝塔分布}

\begin{mydefinition}{贝塔函数与贝塔分布}{beta-function-and-distribution}
设 $a,b>0$，定义\textbf{贝塔函数}
\[
B(a,b)=\int_0^1x^{a-1}(1-x)^{b-1}\,dx.
\]
若随机变量 $X$ 的密度可以取为
\[
f(x)=
\begin{cases}
\dfrac{x^{a-1}(1-x)^{b-1}}{B(a,b)},&0<x<1,\\[0.5em]
0,&\text{其他},
\end{cases}
\]
则称 $X$ 服从参数为 $(a,b)$ 的\textbf{贝塔分布}。
\end{mydefinition}

\textbf{注：}\quad $B(a,b)$ 有限且为正。在 $(0,1/2]$ 上，$(1-x)^{b-1}$ 有界，而 $x^{a-1}$ 可积；在 $[1/2,1)$ 上同理，利用 $(1-x)^{b-1}$ 的可积性。因此上述密度非负，且
\[
\int_{\mathbb R}f(x)\,dx
=\frac1{B(a,b)}\int_0^1x^{a-1}(1-x)^{b-1}\,dx=1.
\]
贝塔分布满足 $P(0<X<1)=1$，可用于描述合格率、市场占有率等比例量。特别地，$B(1,1)=1$，参数 $(1,1)$ 的贝塔分布就是 $U(0,1)$。

\begin{proposition}{贝塔函数的基本性质}{beta-function-properties}
对 $a,b>0$，
\[
B(a,b)=B(b,a),
\qquad
B(a,b)=\frac{\Gamma(a)\Gamma(b)}{\Gamma(a+b)}.
\]
\end{proposition}

\begin{proof}

（1）令 $u=1-x$，则
\[
\begin{aligned}
B(a,b)
&=\int_0^1x^{a-1}(1-x)^{b-1}\,dx\\
&=\int_0^1u^{b-1}(1-u)^{a-1}\,du=B(b,a).
\end{aligned}
\]

（2）由 Gamma 函数定义，将两个积分写为二重积分：
\[
\Gamma(a)\Gamma(b)
=\int_0^\infty\int_0^\infty
u^{a-1}v^{b-1}e^{-(u+v)}\,du\,dv.
\]
在第一象限中令
\[
r=u+v,\qquad x=\frac{u}{u+v},
\]
则 $r>0$、$0<x<1$，反向变换为
\[
u=rx,\qquad v=r(1-x),
\qquad
\left|\frac{\partial(u,v)}{\partial(r,x)}\right|=r.
\]
按二重反常积分换元及非负函数的积分分离，得到
\[
\begin{aligned}
\Gamma(a)\Gamma(b)
&=\int_0^\infty\int_0^1
(rx)^{a-1}\bigl(r(1-x)\bigr)^{b-1}e^{-r}r\,dx\,dr\\
&=\left(\int_0^\infty r^{a+b-1}e^{-r}\,dr\right)
\left(\int_0^1x^{a-1}(1-x)^{b-1}\,dx\right)\\
&=\Gamma(a+b)B(a,b).
\end{aligned}
\]
因为 $\Gamma(a+b)>0$，两边相除即得结论。
\end{proof}

\subsection*{正态分布}

\begin{mydefinition}{正态分布}{normal-distribution}
设 $\mu\in\mathbb R$、$\sigma>0$。若随机变量 $X$ 的密度为
\[
f(x)=\frac1{\sqrt{2\pi}\sigma}
\exp\left(-\frac{(x-\mu)^2}{2\sigma^2}\right),
\qquad x\in\mathbb R,
\]
则称 $X$ 服从参数为 $(\mu,\sigma^2)$ 的\textbf{正态分布}或\textbf{高斯分布}，记为 $X\sim N(\mu,\sigma^2)$。特别地，$N(0,1)$ 称为\textbf{标准正态分布}。
\end{mydefinition}

\begin{proposition}{正态密度的归一化}{normal-density-normalization}
上述 $f$ 非负且在实数轴上的积分为 $1$，因而确实是概率密度。
\end{proposition}

\begin{proof}
非负性由公式直接得到。利用 Gamma 函数部分得到的高斯积分，再作标准化换元：
\[
\int_{-\infty}^{\infty}e^{-y^2/2}\,dy=\sqrt{2\pi}.
\]
此处接续使用该结果。令 $y=(x-\mu)/\sigma$，由于 $\sigma>0$，$dx=\sigma\,dy$，所以
\[
\begin{aligned}
\int_{\mathbb R}f(x)\,dx
&=\frac1{\sqrt{2\pi}\sigma}
\int_{-\infty}^{\infty}e^{-(x-\mu)^2/(2\sigma^2)}\,dx\\
&=\frac1{\sqrt{2\pi}}
\int_{-\infty}^{\infty}e^{-y^2/2}\,dy=1.
\end{aligned}
\]
\end{proof}

\begin{proposition}{正态密度的图像性质}{normal-density-shape}
正态密度关于 $x=\mu$ 对称，在 $\mu$ 处唯一达到最大值 $1/(\sqrt{2\pi}\sigma)$；以 $x$ 轴为水平渐近线，拐点为 $x=\mu\pm\sigma$，图像呈钟形。

固定 $\sigma$ 改变 $\mu$ 时，图像水平平移；固定 $\mu$ 增大 $\sigma$ 时，图像展宽、峰值降低，减小 $\sigma$ 时相反。
\end{proposition}

\begin{proof}

（1）对任意 $h\in\mathbb R$，$f(\mu+h)=f(\mu-h)$，故关于 $x=\mu$ 对称。

（2）求一阶导数，
\[
f'(x)=-\frac{x-\mu}{\sigma^2}f(x).
\]
因为 $f(x)>0$，$x<\mu$ 时导数为正，$x>\mu$ 时为负。因此 $f$ 在 $\mu$ 处唯一达到最大值，代入公式即得其值。

（3）当 $|x|\to\infty$ 时，指数趋于负无穷，故 $f(x)\to0$。再求导，
\[
f''(x)=\frac{(x-\mu)^2-\sigma^2}{\sigma^4}f(x).
\]
导数 $f''$ 在 $|x-\mu|<\sigma$ 时为负，在 $|x-\mu|>\sigma$ 时为正，并在 $\mu\pm\sigma$ 两侧变号，所以这两点为拐点。

（4）把密度写为
\[
f(x)=\frac1{\sigma}\varphi\left(\frac{x-\mu}{\sigma}\right),
\qquad
\varphi(z)=\frac1{\sqrt{2\pi}}e^{-z^2/2}.
\]
可见 $\mu$ 控制水平位置，$\sigma$ 控制水平伸缩及纵向高度，得到最后的图像变化结论。
\end{proof}

\textbf{注：}\quad $\mu$ 称为位置参数，$\sigma$ 通常称为尺度参数，控制曲线的宽窄及峰值高低。

\subsubsection*{标准正态分布与对称性}

标准正态分布的密度和分布函数分别记为
\[
\varphi(x)=\frac1{\sqrt{2\pi}}e^{-x^2/2},
\qquad
\Phi(t)=\int_{-\infty}^t\varphi(x)\,dx.
\]

\begin{proposition}{标准正态分布的对称公式}{standard-normal-symmetry}
对任意 $x\in\mathbb R$，
\[
\Phi(x)+\Phi(-x)=1.
\]
\end{proposition}

\begin{proof}
设 $Z\sim N(0,1)$。利用对称性，由 $\varphi$ 为偶函数，
\[
P(Z\leq-x)=P(Z\geq x).
\]
又因 $P(Z=x)=0$，
\[
\begin{aligned}
\Phi(x)+\Phi(-x)
&=P(Z\leq x)+P(Z\geq x)\\
&=1+P(Z=x)=1.
\end{aligned}
\]
这一论证对正、负 $x$ 及 $x=0$ 均成立。
\end{proof}

\begin{example}[对两侧尾部作符号翻转]
设 $X\sim N(0,1)$，$a>0$，定义
\[
Y=
\begin{cases}
X,&|X|\leq a,\\
-X,&|X|>a.
\end{cases}
\]
证明 $Y\sim N(0,1)$。

先写出分布函数：
\[
\begin{aligned}
F_Y(y)
&=P(X\leq y,\ |X|\leq a)
+P(-X\leq y,\ |X|>a)\\
&=P(X\leq y,\ |X|\leq a)
+P(X\geq-y,\ |X|>a).
\end{aligned}
\]
分三种情况计算。

（1）若 $y>a$，三个相应区间互不相交，
\[
\begin{aligned}
F_Y(y)
&=P(|X|\leq a)+P(-y\leq X<-a)+P(X>a)\\
&=P(X\geq-y)=\Phi(y).
\end{aligned}
\]

（2）若 $-a\leq y\leq a$，利用两侧尾概率相等，
\[
\begin{aligned}
F_Y(y)
&=P(-a\leq X\leq y)+P(X>a)\\
&=P(-a\leq X\leq y)+P(X<-a)\\
&=P(X\leq y)=\Phi(y).
\end{aligned}
\]

（3）若 $y<-a$，中间区间不贡献概率，
\[
F_Y(y)=0+P(X\geq-y)=\Phi(y).
\]
综上 $F_Y=\Phi$，故 $Y\sim N(0,1)$。
\end{example}

\subsubsection*{正态分布的线性变换与标准化}

\begin{mytheorem}{正态分布的线性性质}{normal-affine-transformation}
若 $X\sim N(\mu,\sigma^2)$，$a,b\in\mathbb R$ 且 $a\ne0$，则
\[
aX+b\sim N(a\mu+b,a^2\sigma^2).
\]
\end{mytheorem}

\begin{proof}
令 $Y=aX+b$。用分布函数法，按 $a$ 的符号分别计算。

（1）$a>0$ 时，
\[
\begin{aligned}
F_Y(y)
&=P(aX+b\leq y)\\
&=P\left(X\leq\frac{y-b}{a}\right)
=F_X\left(\frac{y-b}{a}\right).
\end{aligned}
\]
求导得到
\[
f_Y(y)=\frac1a f_X\left(\frac{y-b}{a}\right).
\]

（2）$a<0$ 时，不等号反向。又因 $X$ 的单点概率为 $0$，
\[
\begin{aligned}
F_Y(y)
&=P\left(X\geq\frac{y-b}{a}\right)\\
&=1-F_X\left(\frac{y-b}{a}\right),
\end{aligned}
\]
因此
\[
f_Y(y)=-\frac1a f_X\left(\frac{y-b}{a}\right).
\]

合并两种情形并代入正态密度，
\[
\begin{aligned}
f_Y(y)
&=\frac1{|a|}f_X\left(\frac{y-b}{a}\right)\\
&=\frac1{\sqrt{2\pi}|a|\sigma}
\exp\left(-\frac{(y-a\mu-b)^2}{2a^2\sigma^2}\right).
\end{aligned}
\]
这就是 $N(a\mu+b,a^2\sigma^2)$ 的密度。
\end{proof}

\begin{mycorollary}{正态分布的标准化}{normal-standardization}
若 $X\sim N(\mu,\sigma^2)$，则
\[
\frac{X-\mu}{\sigma}\sim N(0,1).
\]
\end{mycorollary}

\begin{proof}
在线性性质中取 $a=1/\sigma$、$b=-\mu/\sigma$，得到
\[
a\mu+b=0,\qquad a^2\sigma^2=1,
\]
故结论成立。
\end{proof}

\textbf{注：}\quad 对任意 $u<v$，标准化给出计算公式
\[
\begin{aligned}
P(u<X\leq v)
&=P\left(\frac{u-\mu}{\sigma}
<\frac{X-\mu}{\sigma}\leq\frac{v-\mu}{\sigma}\right)\\
&=\Phi\left(\frac{v-\mu}{\sigma}\right)
-\Phi\left(\frac{u-\mu}{\sigma}\right).
\end{aligned}
\]
先标准化，再利用 $\Phi$ 查表或计算；若遇负参数，可用 $\Phi(-x)=1-\Phi(x)$。

\begin{example}[三倍尺度区间的概率]
设 $X\sim N(\mu,\sigma^2)$，则
\[
\begin{aligned}
P(|X-\mu|<3\sigma)
&=P\left(-3<\frac{X-\mu}{\sigma}<3\right)\\
&=\Phi(3)-\Phi(-3)\\
&=2\Phi(3)-1\approx0.9973002.
\end{aligned}
\]
若先取 $\Phi(3)\approx0.9987$，再计算得到 $0.9974$，这是粗查表近似；较精细的概率约为 $99.73\%$。

\textbf{注：}\quad 这称为 $3\sigma$ 原则，可用于理解质量控制中的异常提示。但区间外仍有约 $0.27\%$ 的概率，并非概率为 $0$；一次观测落在区间外，不能单凭这一点就断言随机变量不服从该正态分布。
\end{example}

\begin{example}[由两端比例估计录取分数线]
共有 $10000$ 人报考，成绩大致服从正态分布；$90$ 分以上有 $359$ 人，$60$ 分以下有 $1151$ 人，现从高到低录取 $2500$ 人，估计最低录取分数。

设成绩模型为 $X\sim N(\mu,\sigma^2)$，将人数比例近似为概率。首先，
\[
\begin{aligned}
P(X\geq90)&\approx\frac{359}{10000}=0.0359,\\
\Phi\left(\frac{90-\mu}{\sigma}\right)&\approx0.9641.
\end{aligned}
\]
查表得到
\[
\frac{90-\mu}{\sigma}\approx1.8.
\]
另一方面，
\[
\begin{aligned}
\Phi\left(\frac{60-\mu}{\sigma}\right)
&\approx0.1151,\\
\Phi\left(\frac{\mu-60}{\sigma}\right)
&\approx1-0.1151=0.8849.
\end{aligned}
\]
所以 $(\mu-60)/\sigma\approx1.2$。联立两式，
\[
\begin{cases}
90-\mu\approx1.8\sigma,\\
\mu-60\approx1.2\sigma,
\end{cases}
\qquad
\sigma\approx10,\quad \mu\approx72.
\]

设最低录取分为 $c$。取前 $25\%$，意味着
\[
\begin{aligned}
P(X\geq c)&\approx0.25,\\
\Phi\left(\frac{c-72}{10}\right)&\approx0.75.
\end{aligned}
\]
按粗查表值 $(c-72)/10\approx0.68$，得到 $c\approx78.8$ 分。若用较精细的 $0.75$ 分位数 $0.67449$，在同一 $\mu\approx72,\sigma\approx10$ 的近似下，
\[
c\approx72+10\times0.67449\approx78.74\text{ 分}.
\]
$78.8$ 分是粗查表结果，不是精确分位数。
\end{example}

\section{随机变量函数的分布}
\label{sec:distributions-of-functions}
\subsection*{随机变量的函数与离散型随机变量函数的分布列}

设 $X$ 定义在概率空间 $(\Omega,\mathcal F,P)$ 上，$g:\mathbb R\to\mathbb R$。复合映射
\[
Y(\omega)=g(X(\omega))
\]
确实是从 $\Omega$ 到 $\mathbb R$ 的映射，但要称它为随机变量，还需要可测性。对于任意博雷尔集 $B$，有
\[
Y^{-1}(B)=X^{-1}\bigl(g^{-1}(B)\bigr).
\]
因此，若 $g^{-1}(B)$ 总是博雷尔集，就可由 $X$ 的可测性保证 $Y^{-1}(B)\in\mathcal F$。这是对任意随机变量 $X$ 都适用的充分条件；对于某个固定的 $X$，不必声称它是必要条件。

\begin{mydefinition}{博雷尔函数}{borel-function}
若函数 $g:\mathbb R\to\mathbb R$ 满足
\[
g^{-1}(B)\in\mathcal B(\mathbb R),
\qquad \forall B\in\mathcal B(\mathbb R),
\]
则称 $g$ 为\textbf{博雷尔函数}，也称 $(\mathbb R,\mathcal B(\mathbb R))$ 上的可测函数。
\end{mydefinition}

若离散型随机变量 $X$ 的不同可能取值为 $x_k$，记 $p_k=P(X=x_k)$。将所有使 $g(x_k)$ 相同的项合并，即得
\[
P(Y=y)=\sum_{\{k:g(x_k)=y\}}p_k.
\]
理由是 $\{Y=y\}$ 在 $X$ 的这些可能取值上，分解为互不相交的事件 $\{X=x_k\}$；这些取值以外的事件概率为零，可以忽略。

\begin{example}[离散型随机变量的多项式变换]
已知 $X$ 的分布列为
\[
\begin{array}{c|ccccc}
X&-2&-1&0&1&2\\ \hline
P&0.2&0.1&0.1&0.3&0.3
\end{array},
\]
求 $Y=X^2+X$ 的分布列。

逐项计算并保留原来的概率：
\[
\begin{array}{c|ccccc}
X&-2&-1&0&1&2\\ \hline
Y=X^2+X&2&0&0&2&6\\
P&0.2&0.1&0.1&0.3&0.3
\end{array}.
\]
$Y=0$ 对应 $X=-1,0$，$Y=2$ 对应 $X=-2,1$，$Y=6$ 对应 $X=2$。合并后得到
\[
\begin{array}{c|ccc}
Y&0&2&6\\ \hline
P&0.2&0.5&0.3
\end{array}.
\]
三个概率之和为 $1$。
\end{example}

\subsection*{求随机变量函数密度的一般方法}

\begin{example}[一般连续型随机变量的线性变换]
设 $X$ 的密度为 $f_X$，求 $Y=aX+b$（$a\ne0$）的密度。

先按 $a$ 的符号求分布函数。$a>0$ 时，
\[
\begin{aligned}
F_Y(t)
&=P\left(X\leq\frac{t-b}{a}\right)\\
&=F_X\left(\frac{t-b}{a}\right).
\end{aligned}
\]
$a<0$ 时，利用 $X$ 无单点质量，
\[
\begin{aligned}
F_Y(t)
&=P\left(X\geq\frac{t-b}{a}\right)\\
&=1-F_X\left(\frac{t-b}{a}\right).
\end{aligned}
\]
分别对上式求导，合并得到一个密度版本
\[
\boxed{\quad f_Y(t)=\frac1{|a|}f_X\left(\frac{t-b}{a}\right).\quad}
\]
\end{example}

一般方法是“先分布函数，后密度”。对使 $g(X)$ 为随机变量的可测函数 $g$，令
\[
D(t)=\{x\in\mathbb R:g(x)\leq t\}.
\]
先解出这个取值集合，再计算
\[
\begin{aligned}
F_Y(t)&=P(g(X)\leq t)\\
&=P(X\in D(t))=\int_{D(t)}f_X(x)\,dx.
\end{aligned}
\]
若已确认 $Y$ 具有密度，则由 $F_Y$ 几乎处处求导可得到密度版本。
\subsection*{构造给定分布函数的随机变量}

\begin{mydefinition}{分布函数的广义逆}{generalized-inverse-cdf}
设 $G$ 单调不减、右连续，且
\[
\lim_{x\to-\infty}G(x)=0,\qquad
\lim_{x\to+\infty}G(x)=1.
\]
对 $p\in(0,1)$，定义其\textbf{广义逆}
\[
G^{-1}(p)=\inf\{x\in\mathbb R:G(x)\geq p\}.
\]
这里不要求 $G$ 严格单调，也不要求它有通常意义的反函数。
\end{mydefinition}

\textbf{注：}\quad 广义逆在 $p\in(0,1)$ 时是有限实数。由两端极限，可以找到 $L,R$ 使 $G(L)<p<G(R)$。因此集合
\[
A_p=\{x:G(x)\geq p\}
\]
非空，且其中任意点都大于 $L$，所以其下确界存在且有限。

\begin{mylemma}{广义逆与分布函数的不等式}{generalized-inverse-equivalence}
对任意 $p\in(0,1)$ 和 $t\in\mathbb R$，
\[
G^{-1}(p)\leq t\quad\Longleftrightarrow\quad p\leq G(t).
\]
\end{mylemma}

\begin{proof}
记 $q=G^{-1}(p)=\inf A_p$。分别证明两个方向，并在等号边界使用右连续性。

\textbf{（$\Leftarrow$）}\quad 若 $p\leq G(t)$，则 $t\in A_p$，从而
\[
q=\inf A_p\leq t.
\]

\textbf{（$\Rightarrow$）}\quad 设 $q\leq t$，分两种情况。

（1）若 $t>q$，由下确界定义，存在 $x\in A_p$ 使 $x<t$；否则 $t$ 也是 $A_p$ 的下界，与 $q<t$ 矛盾。由 $G$ 单调不减，
\[
G(t)\geq G(x)\geq p.
\]

（2）若 $t=q$，令 $t_n=t+1/n$，则 $t_n>q$ 且 $t_n\downarrow t$。由（1），每个 $n$ 都满足 $G(t_n)\geq p$。再由右连续性，
\[
G(t)=\lim_{n\to\infty}G(t_n)\geq p.
\]
两个方向均得证。
\end{proof}

\begin{mytheorem}{由均匀分布构造给定分布}{inverse-transform-construction}
设 $X\sim U(0,1)$，$G$ 及其广义逆满足前述定义。令
\[
Y=G^{-1}(X),
\]
则 $Y$ 的分布函数就是 $G$。
\end{mytheorem}

\begin{proof}
首先处理定义域。因为 $P(0<X<1)=1$，广义逆在这里为有限实数；在零概率事件 $\{X\notin(0,1)\}$ 上，可将 $Y$ 另定义为 $0$。

对任意 $y\in\mathbb R$，在 $\{0<X<1\}$ 上，由前述引理，
\[
\{G^{-1}(X)\leq y\}=\{X\leq G(y)\}.
\]
记 $E=\{X\notin(0,1)\}$。由 $Y$ 在 $E$ 上定义为 $0$，
\[
\{Y\leq y\}
=\bigl(E^c\cap\{X\leq G(y)\}\bigr)
\cup
\begin{cases}
E,&y\geq0,\\
\varnothing,&y<0.
\end{cases}
\]
右侧由可测事件取交、取并构成，所以 $Y$ 是随机变量；这里不需要假定零概率集的任意子集都可测。

$G$ 单调不减且两端极限为 $0$ 和 $1$，所以 $0\leq G(y)\leq1$。由均匀分布，
\begin{align*}
F_Y(y)
&=P(Y\leq y)\\
&=P\bigl(X\leq G(y)\bigr)\\
&=\int_0^{G(y)}1\,dx
=G(y).
\end{align*}
这一计算也包括 $G(y)=0$ 或 $1$ 的情形。
\end{proof}

\textbf{注：}\quad 因此，单调不减、右连续且两端极限为 $0$、$1$，不仅是分布函数的必要性质，也足以构造一个以它为分布函数的随机变量。可以从 $(0,1)$ 上的均匀概率空间出发，取恒等随机变量 $X(\omega)=\omega$，再令 $Y=G^{-1}(X)$。这一结论不要求 $G$ 连续或严格递增。

\subsection*{求随机变量函数密度的简单方法}

\begin{mytheorem}{分支反函数的密度变换公式}{inverse-branches-density}
设 $X$ 有密度 $f_X$，$g$ 为博雷尔函数，$Y=g(X)$。设存在有限个或可数个互不相交的开区间 $I_i$，满足
\[
P\left(X\in\bigcup_i I_i\right)=1.
\]
在每个 $I_i$ 上，$g$ 连续、严格单调，像 $J_i=g(I_i)$ 为开区间；其反函数 $h_i:J_i\to I_i$ 连续可微。这里 $h_i$ 称为第 $i$ 个\textbf{反函数分支}，即
\[
g(h_i(y))=y,\qquad y\in J_i.
\]
则 $Y$ 有一个密度版本
\[
f_Y(y)=\sum_{\{i:y\in J_i\}}f_X(h_i(y))\,|h_i'(y)|.
\]
没有分支贡献时，该和取 $0$。非负级数在几乎处处有限；在可能出现的零测异常点上，可另赋有限值而不改变分布。
\end{mytheorem}

\begin{example}[用反函数公式求线性变换的密度]
设 $Y=aX+b$，$a\ne0$。变换在整个实轴上只有一个反函数分支
\[
h(y)=\frac{y-b}{a},\qquad |h'(y)|=\frac1{|a|}.
\]
所以
\[
f_Y(y)=\frac1{|a|}f_X\left(\frac{y-b}{a}\right).
\]
这与分布函数法的结果一致。
\end{example}

\begin{example}[平方变换与不同分支的取值范围]
设随机变量 $X$ 的密度为
\[
f_X(x)=
\begin{cases}
\dfrac14|x|,&-2\leq x<0,\\
x,&0\leq x\leq1,\\
0,&\text{其他}.
\end{cases}
\]
求 $Y=X^2$ 的密度。

负半轴的概率为 $\int_{-2}^0|x|/4\,dx=1/2$，正半轴的概率为 $\int_0^1x\,dx=1/2$，所以给定密度归一化。由 $-2\leq X\leq1$ 以概率 $1$ 成立，$Y$ 可限制在 $[0,4]$。

对 $y>0$，两个反函数分支为 $h_1(y)=\sqrt y$ 和 $h_2(y)=-\sqrt y$，它们的导数绝对值都为 $1/(2\sqrt y)$，因此
\[
f_Y(y)=\frac{f_X(\sqrt y)+f_X(-\sqrt y)}{2\sqrt y}.
\]
当 $0<y<1$ 时，正、负两个原像均落在原密度的非零区间内，故
\[
f_Y(y)=\frac{\sqrt y+\sqrt y/4}{2\sqrt y}=\frac58.
\]
当 $1<y<4$ 时，正原像超过 $1$，只有负分支贡献，故
\[
f_Y(y)=\frac{0+\sqrt y/4}{2\sqrt y}=\frac18.
\]
当 $y<0$ 或 $y>4$ 时，概率密度可取 $0$。在 $y=0$ 处不能代入含 $1/\sqrt y$ 的公式；但
\[
P(Y=0)=P(X=0)=0,
\]
可另指定该点的密度值。取密度版本为
\[
f_Y(y)=
\begin{cases}
\dfrac58,&0\leq y\leq1,\\
\dfrac18,&1<y\leq4,\\
0,&\text{其他}.
\end{cases}
\]
检查全积分：
\[
\int_{\mathbb R}f_Y(y)\,dy=\frac58+\frac18(4-1)=1.
\]
\end{example}

\begin{example}[均匀随机变量的余弦变换]
设 $X\sim U(0,6\pi)$，求 $Y=\cos X$ 的密度。

$Y\in[-1,1]$。对 $-1<y<1$，令 $\theta=\arccos y\in(0,\pi)$。区间 $(0,6\pi)$ 包含三个完整周期，每个周期上余弦先减后增，所以所有原像恰为六个互不相同的点
\[
h_{k,1}(y)=\arccos y+2k\pi,\qquad
h_{k,2}(y)=2\pi-\arccos y+2k\pi,
\quad k=0,1,2.
\]
六个反函数分支的导数绝对值均为 $1/\sqrt{1-y^2}$，原密度均为 $1/(6\pi)$，于是
\begin{align*}
f_Y(y)
&=\sum_{k=0}^2
\left[
\frac1{6\pi\sqrt{1-y^2}}+
\frac1{6\pi\sqrt{1-y^2}}
\right]\\
&=\frac1{\pi\sqrt{1-y^2}},\qquad -1<y<1.
\end{align*}
在 $y=\pm1$ 处，原像为有限个点，概率为零；分支导数公式在端点不适用，可将密度另取为 $0$。所以一个完整版本为
\[
f_Y(y)=
\begin{cases}
\dfrac1{\pi\sqrt{1-y^2}},&-1<y<1,\\
0,&\text{其他}.
\end{cases}
\]
虽然密度在两端附近无界，但积分收敛，并且
\[
\int_{-1}^1\frac{dy}{\pi\sqrt{1-y^2}}
=\frac1\pi\left[\arcsin y\right]_{-1}^1=1.
\]
\end{example}

\begin{example}[倒数的小数部分与可数个反函数分支]
设 $X\sim U(0,1)$，令 $Y$ 为 $1/X$ 的小数部分。记 $\lfloor u\rfloor$ 为不超过 $u$ 的最大整数，即向下取整，因而
\[
Y=\frac1X-\left\lfloor\frac1X\right\rfloor\in[0,1).
\]
在零概率事件 $\{X\notin(0,1)\}$ 上，可另定义 $Y=0$。

对 $0<y<1$，若 $Y=y$，则唯一的整数部分 $k=\lfloor1/X\rfloor$ 满足 $k\geq1$，且
\[
\frac1X=k+y
\quad\Longleftrightarrow\quad X=h_k(y)=\frac1{k+y}.
\]
反过来，每个 $k\geq1$ 给出的 $h_k(y)\in(0,1)$ 都是有效原像。各分支对应互不相交的区间
\[
I_k=\left(\frac1{k+1},\frac1k\right).
\]
这些区间只遗漏可数个点 $1/k$，其总概率为零，满足分支公式的全概率条件。由于
\[
|h_k'(y)|=\frac1{(k+y)^2},
\]
得到一个密度版本
\[
f_Y(y)=
\begin{cases}
\displaystyle\sum_{k=1}^{\infty}\frac1{(k+y)^2},&0\leq y<1,\\
0,&\text{其他}.
\end{cases}
\]
这里在 $y=0$ 处的值是另选的密度值；该点的原像是可数集，$P(Y=0)=0$。对 $0\leq y<1$，各项不超过 $1/k^2$，所以级数收敛。

利用非负求和与积分交换，再作望远镜求和，核验
\begin{align*}
\int_0^1 f_Y(y)\,dy
&=\sum_{k=1}^{\infty}\int_0^1\frac{dy}{(k+y)^2}\\
&=\sum_{k=1}^{\infty}\left(\frac1k-\frac1{k+1}\right)
=1.
\end{align*}
\end{example}
